%\setcounter{framenumber}{30}
\renewcommand{\mydate}{2025年03月27号\ 第十次课}

\makeatletter
\ifodd\value{page}
\else
\begin{frame}[plain]
\end{frame}
\fi
\makeatother


\section{矩阵(二)}


\title{第五章 \quad 矩阵(二)}
\author{}
\date{}

\frame{\titlepage}


\subsection{初等变换}


\subsubsection{初等矩阵}
\begin{frame}\frametitle{初等矩阵}
\begin{itemize}
\item 初等行变换
\pp
\begin{itemize}
\item 交换两行
\item 某行乘以一个非零常数
\item 某行乘以一个常数加到另一行上
\end{itemize}
\pp
\item 初等列变换
\pp
\begin{itemize}
\item 交换两列
\item 某列乘以一个非零常数
\item 某列乘以一个常数加到另一列上
\end{itemize}
\end{itemize}
\pp

\blue{初等变换}(6个) = \blue{初等行变换}(3个) + \blue{初等列变换}(3个)

\bigskip

\blue{问题:} 对单位矩阵做初等变换,我们会得到什么矩阵?
\end{frame}








\begin{frame}\frametitle{初等矩阵}
\begin{definition}[初等矩阵]
\begin{enumerate}
\item 交换单位阵的第$i,j$行(或$i,j$列)得
\[\tiny \blue{S_{ij}} := \begin{pmatrix}
1\quad&\qquad&\qquad&\qquad&\qquad&\qquad&\qquad\\
  &\ddots&      &      &      &      &     \\
  &      &  0   &      &  1   &      &     \\
  &      &      &\ddots&      &      &     \\
  &      &  1   &      &  0   &      &     \\
  &      &      &      &      &\ddots&     \\
  &      &      &      &      &      &  1  \\

\end{pmatrix}\]
\pp
\item 将单位阵第$i$行(或$i$列)乘以非零常数$\lambda$得
\[\tiny \blue{D_i(\lambda)} :=\begin{pmatrix}
1\quad&\qquad&\qquad&\qquad&\qquad&\qquad&\qquad\\
&\ddots&      &       &      &      &     \\
&      &  1   &       &      &      &     \\
&      &      &\lambda&      &      &     \\
&      &      &       &  1   &      &     \\
&      &      &       &      &\ddots&     \\
&      &      &       &      &      &  1  \\
\end{pmatrix}\]
\end{enumerate}
\end{definition}
\end{frame}





\begin{frame}\frametitle{初等矩阵}
\begin{definition}[初等矩阵]
\begin{enumerate}\setcounter{enumi}{2}
\item 将单位阵的第$j$行的$\lambda$倍加到第$i$行(或第$i$列的$\lambda$倍加到第$j$列)得
\[\tiny \blue{T_{ij}(\lambda)} := \begin{pmatrix}
1\quad&\qquad&\qquad&\qquad&\qquad&\qquad&\qquad\\
&\ddots&      &      &       &      &     \\
&      &  1   &      &\lambda&      &     \\
&      &      &\ddots&       &      &     \\
&      &      &      &  1    &      &     \\
&      &      &      &       &\ddots&     \\
&      &      &      &       &      &  1  \\
\end{pmatrix}\]
\end{enumerate}
\pp
我们称上面定义的三类矩阵为\blue{初等矩阵}.
\end{definition}

\end{frame}









% ========== 2023年04月06号 第10次课 初等矩阵 ==========
%\frame{\titlepage}





\subsubsection{初等矩阵基本性质}
\begin{frame}\frametitle{初等矩阵基本性质}

\begin{theorem}
\begin{enumerate}
\item 对矩阵做初等\blue{行}变换相当于矩阵\blue{左}乘一个相应的初等矩阵.
\pp
\item 对矩阵做初等\blue{列}变换相当于矩阵\blue{右}乘一个相应的初等矩阵.
\end{enumerate}
\end{theorem}
\pp
证明思路:将矩阵写成由向量组成的分块矩阵,然后直接验证.
\pp

\begin{proposition}[初等矩阵总是可逆的，且其逆仍然是初等矩阵]
\begin{enumerate}
\item $S_{ij}$对称且$S_{ij}^{-1}=S_{ij}$; \pp
\item $D_i(\lambda)$为对角阵且$D_i(\lambda)^{-1}=D_i(\lambda^{-1})$; \pp
\item $T_{ij}(\lambda)$为三角阵,且$T_{ij}(\lambda)^{-1}=T_{ij}(-\lambda)$;
\end{enumerate}
\end{proposition}
\pp
特别地，由于可逆矩阵的乘积仍然可逆，我们有任意有限多个初等矩阵的乘积也是可逆的.
\end{frame}



\subsection{初等矩阵的应用}

\subsubsection{高斯消元法与可逆阵的$LU$分解.}
\begin{frame}\frametitle{高斯消元法与可逆阵的$LU$分解.}
由高斯消元法,对任意矩阵$A$可通过初等行变换将其变为阶梯矩阵.
\pp
用初等矩阵,这过程可描述为: 存在初等方阵$P_1,P_2,\cdots,P_s$使得
\[P_s\cdots P_2P_1A =: J\]
为阶梯矩阵.
\pp
\bigskip
\begin{itemize}
\item 若 $A$ 可逆，则$J$为可逆上三角矩阵.
\pp
\bigskip
\item 若 $P_1$, $P_2$, $\cdots$, $P_s$ 仅出现下三角初等矩阵 $T_{ij}(\lambda)$($j<i$). 则记 $L:=P_1^{-1}P_2^{-1}\cdots P_s^{-1}$ 以及 $U:=J$, 称 $A=LU$\footnote{注：$L$的$(i,j)$元素正好为依次将$A$的$(i,j)$元素化为$0$所需减去的$(j,j)$元素的倍数.}为$A$的\blue{$LU$分解}.
\pp


\pp
\bigskip



\item 若可逆阵$A$的$LU$分解存在，则将解方程组 $Ax=b$可转化为两个简单方程组$Ly=b$, $Ux=y$ (这两个方程都可通过带入法求解).
\end{itemize}
\end{frame}




\subsubsection{用LU分解求解线性方程组}
\begin{frame}\frametitle{用LU分解求解线性方程组}
\begin{example}[用LU分解求解线性方程组]
\[\begin{cases}
2x_1 + 3x_2 + 4x_3 = 1 \\
4x_1 + 11x_2 + 14x_3 = 3 \\
6x_1 + 19x_2 + 26x_3 = 5.
\end{cases}\]
\end{example}
\begin{itemize}
\item 第一步：对系数矩阵进行LU分解得（Gauss消元法）：
\[L = \begin{pmatrix}
1 & 0 & 0 \\
2 & 1 & 0 \\
3 & 2 & 1
\end{pmatrix}, \quad
U = \begin{pmatrix}
2 & 3 & 4 \\
0 & 5 & 6 \\
0 & 0 & 2
\end{pmatrix}.\]
\pp

\item 第二步：求解 $Ly = b$. 这一步通过自上而下带入完成.
\[\begin{pmatrix}
1 & 0 & 0 \\
2 & 1 & 0 \\
3 & 2 & 1
\end{pmatrix}
\begin{pmatrix}y_1 \\y_2 \\y_3\end{pmatrix}
=
\begin{pmatrix}1 \\3 \\5\end{pmatrix}
\Rightarrow y = \begin{pmatrix} 1 \\ 1 \\ 0 \end{pmatrix}.\]
\pp

\item 第三步：求解 $Ux = y$. 通过回代（back substitution）
\[\begin{pmatrix}
2 & 3 & 4 \\
0 & 5 & 6 \\
0 & 0 & 2
\end{pmatrix}
\begin{pmatrix}x_1 \\x_2 \\x_3\end{pmatrix}
=
\begin{pmatrix}1 \\1 \\0\end{pmatrix}
\Rightarrow x = \begin{pmatrix} \frac{1}{5} \\ \frac{1}{5} \\ 0 \end{pmatrix}.\]
\end{itemize}
\end{frame}




\subsubsection{初等矩阵与高斯消元法}
\begin{frame}\frametitle{初等矩阵与高斯消元法}
\begin{theorem}对于任意矩阵$A\in F^{m\times n}$,存在$m$阶初等方阵$P_1,P_2,\cdots,P_s$和$n$阶初等方阵$Q_1,Q_2,\cdots,Q_t$，以及非负整数$r$，使得
\[P_s\cdots P_2P_1AQ_1Q_2\cdots Q_t=\begin{pmatrix}I_r&0\\0&0\end{pmatrix}.\]
\end{theorem}
\pp
\begin{corollary}
\begin{enumerate}
\item 对任意矩阵$A\in \mathbb F^{m\times n}$,存在$m$阶可逆阵$P$和$n$阶可逆阵$Q$使得
\[PAQ=\begin{pmatrix}I_r&0\\0&0\end{pmatrix}.\]
\pp
\item 设$A$为$n$阶方阵.则$A$可逆当且仅当$A$可分解为初等方阵的乘积.
\end{enumerate}
\end{corollary}

\end{frame}










\renewcommand{\mydate}{2025年04月01号\ 第十一次课}












\subsubsection{初等变换求逆}
\begin{frame}\frametitle{初等变换求逆}

\begin{corollary}设$A$为$n$阶可逆方阵.则
\pp
\begin{enumerate}
\item 可对$A$做一系列初等行变换变为最简形式 $I$.
\pp
\item 可对$A$做一系列初等列变换变为最简形式 $I$.
\end{enumerate}
\end{corollary}
\pp
\bigskip

\blue{问题:} 如何求逆矩阵?

\pp
\bigskip


\begin{itemize}
\item \blue{算法原理:} 通过初等变换我们可以找到初等矩阵 $P_1,P_2,\cdots,P_s$使得
$P_s\cdots P_2P_1A=I$. \pp 从而,我们有
\[A^{-1} =  P_s\cdots P_2P_1.\]
 \pp

\item \blue{具体实现过程:} 对矩阵$(A,I)$做\blue{行}初等变换,将第一个子矩阵变为单位阵.则第二个子矩阵自动变为$A$的逆. \pp
\[(A,I)\xrightarrow{\text{行初等变换}} P_s\cdots P_2P_1(A,I)=(I,A^{-1}).\]
\end{itemize}
\end{frame}







\begin{frame}\frametitle{初等变换求逆}
\begin{example}
求矩阵$\tiny\begin{pmatrix}2&1&0\\1&2&1\\0&1&2\end{pmatrix}$的逆. \pp \qquad \qquad \qquad  \yang{$\tiny\frac14\begin{pmatrix}3&-2&1\\-2&4&-2\\1&-2&3\end{pmatrix}$.}
\end{example}
\pp

思考: 若$A$可逆,如何求解 $AX=B$? \pp
\[(A,B)\xrightarrow{\text{行初等变换}} P_s\cdots P_2P_1(A,B)=(I,A^{-1}B).\]
\pp

\begin{example}
若 $\tiny\begin{pmatrix}2&1&0\\1&2&1\\0&1&2\end{pmatrix}X = \begin{pmatrix} 3& 1\\ 4& 0\\ 3& 1\\ \end{pmatrix}$, 求 $X$. \pp \qquad \qquad \qquad   \yang{$\tiny\begin{pmatrix}1&1\\1&-1\\1&1\\\end{pmatrix}$.}
\end{example}

思考: 若$A$可逆,如何求解 $XA=B$? \pp
\[\begin{pmatrix}A\\B\\\end{pmatrix} \xrightarrow{\text{列初等变换}} \begin{pmatrix}A\\B\\\end{pmatrix}Q_1Q_2\cdots Q_t=\begin{pmatrix}I\\BA^{-1}\\\end{pmatrix}.\]
\pp

\begin{example}
若 $\tiny X\begin{pmatrix}2&1&0\\1&2&1\\0&1&2\end{pmatrix} = \begin{pmatrix} 3&4&3\\1&0&1\\\end{pmatrix}$, 求 $X$. \pp \qquad \qquad \qquad   \yang{$\tiny\begin{pmatrix}1&1&1\\1&-1&1\\\end{pmatrix}$.}
\end{example}

\end{frame}


\subsection{分块矩阵的初等变换及应用}

\subsubsection{分块矩阵的行列变换}
\begin{frame}\frametitle{分块矩阵的行列变换}
设$A$为$m\times n$矩阵. 我们将其按$m=m_1+m_2+\cdots+m_r$和$n=n_1+n_2+\cdots+n_s$的拆分方式,得到如下分块形式
\[ A = \left(\begin{array}{cccc}
A_{11} & A_{12} & \cdots & A_{1 s}\\
A_{21} & A_{22} & \cdots & A_{2 s}\\
\vdots & \vdots & \ddots & \vdots\\
A_{r 1} & A_{r 2} & \cdots & A_{rs}
\end{array}\right) = (A_{ij})_{r \times s}.\]
\pp
类似于矩阵的初等变换,我们可以对分块矩阵$A$做相似的操作:\pp
\begin{enumerate}
\item 交换分块矩阵的第$i,j$行(或$i,j$列). \pp
\item 将分块矩阵的第$i$行左乘一个$m_i$阶(或,第$i$列右乘一个$n_i$阶)可逆矩阵. \pp
\item 将分块矩阵的第$i$行左乘一个$m_j\times m_i$矩阵加到第$j$行(或,第$i$列右乘一个$n_i\times n_j$矩阵加到第$j$列).
\end{enumerate}
\end{frame}





\subsubsection{初等分块矩阵}
\begin{frame}\frametitle{初等分块矩阵}
\begin{example}若$A$可逆,则
\begin{equation*} \footnotesize
\begin{split}
&
\begin{pmatrix}I&0\\-CA^{-1}&I\end{pmatrix} \begin{pmatrix}A&B\\C&D\end{pmatrix} =\begin{pmatrix}A&B\\0&D-CA^{-1}B\end{pmatrix}\\ \pp
&\\
&
\begin{pmatrix}A&B\\C&D\end{pmatrix} \begin{pmatrix}I&-A^{-1}B\\0&I\end{pmatrix} =\begin{pmatrix}A&0\\C&D-CA^{-1}B\end{pmatrix}\\ \pp
&\\
&
\begin{pmatrix}I&0\\-CA^{-1}&I\end{pmatrix}\begin{pmatrix}A&B\\C&D\end{pmatrix} \begin{pmatrix}I&-A^{-1}B\\0&I\end{pmatrix} =\begin{pmatrix}A&0\\0&D-CA^{-1}B\end{pmatrix}\\
\end{split}
\end{equation*}
\end{example}
\end{frame}





\subsubsection{分块矩阵求逆}
\begin{frame}\frametitle{分块矩阵求逆}

\begin{example}设$A,B,I$为$n$阶方阵满足$BA=0$.则
\[\small \begin{pmatrix}I&A\\B&I\end{pmatrix}^{-1}=\begin{pmatrix}I+AB&-A\\-B&I\end{pmatrix}\]
\end{example}
\pp
\begin{equation*}
\begin{split}
\begin{pmatrix}I&A&I&0\\B&I&0&I\end{pmatrix}
& \xrightarrow[\text{第一行左乘$B$}]{\text{第二行减去}}
\begin{pmatrix}I&A&I&0\\0&I&-B&I\end{pmatrix} \\ \pp
&\\
& \xrightarrow[\text{第二行左乘$A$}]{\text{第一行减去}}
\begin{pmatrix}I&0&I+AB&-A\\0&I&-B&I\end{pmatrix} \\
\end{split}
\end{equation*}
\end{frame}





\subsection{秩与相抵}


\subsubsection{秩与相抵}
\begin{frame}\frametitle{秩与相抵}
\begin{definition}[相抵] 称两矩阵$A,B$ \blue{相抵},若存在可逆矩阵$P$和$Q$使得
\[B=PAQ.\]
\end{definition}
\pp
注:因可逆矩阵为方阵,因此要使$B=PAQ$成立, $B$和$A$的行列数必须一致.
\pp

\begin{proposition}[相抵的基本性质]
\begin{enumerate}
\item  任意矩阵与自身相抵; \pp
\item  若$A$与$B$相抵,则$B$与$A$也相抵; \pp
\item  若$A$与$B$相抵且$B$与$C$相抵,则$A$与$C$相抵.
\end{enumerate}
\end{proposition}
满足上述三条性质的关系称为\blue{等价关系}。换句话说，上述性质表明相抵为等价关系.
\end{frame}





\subsubsection{等价关系在集合论中的严谨定义}
\begin{frame}\frametitle{等价关系在集合论中的严谨定义}
\begin{definition}[等价关系严谨定义]
给定集合$X$,我们称$X\times X$的一个子集为$X$上的一个\blue{关系}. \pp 设 $R$ 为 $X$ 上的一个关系, 即
\[R\subseteq X\times X.\] \pp
我们称 $R$是\blue{等价关系}, 若 $R$ 满足如下三条性质
\begin{itemize}
\item (\blue{自反性}) 对任意 $x\in X$, 我们有 $(x,x)\in R$;
\item (\blue{对称性}) 对任意 $x,y\in X$, $(x,y)\in R$ 当且仅当 $(y,x)\in R$;
\item (\blue{传递性}) 对任意 $x,y,z\in X$, 若 $(x,y)\in R$, $(y,z)\in R$ 则 $(x,z)\in R$.
\end{itemize}
\pp
此时,对任意$x,y\in X$, 若$(x,y)\in R$, 则称\blue{$x$等价于$y$}; 否则,称\blue{$x$不等价于$y$}.
\end{definition}

\begin{example}
例子：老乡；同号；同余等价关系；

反例：朋友；父子；异号
\end{example}

\end{frame}




\subsubsection{集合上等价关系与拆分}
\begin{frame}\frametitle{集合上等价关系与拆分}
对等价关系,我们有如下通俗理解:
\begin{theorem}
给定集合某上的一个等价关系等价于给定这集合上的一个拆分.
\end{theorem}
\pp 证明思路：
\begin{itemize}
\item 给定$X$上一个等价关系 $R$. 我们称$X$的子集
\[[x]:=\{y\in X\mid (x,y)\in R\} \subseteq X\]
为$x$在等价关系$R$下的\blue{等价类}. \pp
\begin{itemize}
\item 对任意$x,y\in X$, $[x]\cap[y]\neq\emptyset$ 当且仅当 $[x]=[y]$. \pp
\item 集合$X$ 为一些等价类的无交并. \pp
\end{itemize}
\item 反之, 给定集合$X$的一个拆分 $X$. 我们如下定义一个等价关系$R$: 对于任意$x,y\in X$, $(x,y)\in R$当且仅当$x,y$落入同一个拆分子集中.
\end{itemize}
\end{frame}



\subsubsection{相抵等价关系中的两个基本问题}
\begin{frame}\frametitle{秩与相抵}

\pp
按照相抵关系,我们将全体$m\times n$矩阵分成若干个两两不交的等价类. \pp

\blue{基本问题:}
\begin{enumerate}
\item 如何判定两个矩阵是否相抵? \pp
\item 在每个等价类中有没有最简单的表达形式?
\end{enumerate}

\end{frame}




\begin{frame}\frametitle{秩与相抵}
\begin{theorem}
对于任意矩阵$A\in F^{m\times n}$,存在$m$阶可逆矩阵$P$和$n$阶可逆矩阵$Q$使得
\[PAQ=\begin{pmatrix}I_r&0\\0&0\end{pmatrix}.\]
其中非负整数$r$\blue{不依赖}于$P$和$Q$的选取.
\end{theorem}
\pp
\bigskip

\yang{假若不唯一，不妨设$s>r$ 且 $P\begin{pmatrix}I_r&0\\0&0\end{pmatrix}Q = \begin{pmatrix}I_s&0\\0&0\end{pmatrix}$. 对$P$和$Q$进行分块：
\[P = \begin{pmatrix}(P_{11})_{s\times r}
& (P_{12})_{s\times(m-r)}\\
(P_{21})_{(m-s)\times r}
& (P_{22})_{(m-s)\times(m-r)}\\ \end{pmatrix}_{m\times m}\]
\[Q = \begin{pmatrix}(Q_{11})_{r\times s}
& (Q_{12})_{r\times(n-s)}\\
(Q_{21})_{(n-r)\times s}
& (Q_{22})_{(n-r)\times(n-s)}\\ \end{pmatrix}_{n\times n}.\]
则
$P_{11}Q_{11}=I_s$(矛盾！).}
\end{frame}









% ========== 2023年04月11号 第11次课 秩与相抵 ==========
%\frame{\titlepage}





\subsubsection{相抵判定}
\begin{frame}\frametitle{相抵判定}

\begin{definition}
定理中的$\begin{pmatrix}I_r&0\\0&0\end{pmatrix}$称为$A$的\blue{相抵标准形}.\pp 非负整数$r$称为$A$的\blue{秩},\pp 记为\blue{$\rank(A)$}或\blue{$r(A)$}.\pp 若$r=m$,则称$A$为\blue{行满秩}.\pp
若$r=n$,则称$A$为\blue{列满秩}.
\end{definition}

\begin{corollary}[相抵判定]
给定两矩阵$A,B\in F^{m\times n}$,则$A$与$B$相抵当且仅当$\rank(A)=\rank(B)$.
\end{corollary}
\pp
\begin{corollary}若$P,Q$可逆,则$\rank(PAQ)=\rank(A)$.
\end{corollary}
\pp
\begin{example}
\[\rank(A^T)=\rank(A).\]
\pp
\[\rank\begin{pmatrix}A&0\\0&B\end{pmatrix}=\rank(A)+\rank(B).\]
\end{example}
\end{frame}

\renewcommand{\mydate}{2025年04月03号\ 第十二次课}



\subsubsection{通过子式定义秩}
\begin{frame}\frametitle{通过子式定义秩}
\begin{proposition}
设$A\in F^{m\times n}$, $P$, $Q$分别为$m$, $n$阶初等方阵.若$A$的$k$阶子式全为零,则$PA$和$AQ$的所有$k$阶子式全为零.
\end{proposition}
\pp

\begin{proposition}[通过子式定义秩]
若矩阵$A$至少有一个$r$阶非零子式,且$A$的所有$r+1$阶子式都为零,则$\rank(A)=r$.
\end{proposition}
\end{frame}




\subsection{相关例题}

\begin{frame}\frametitle{例题}
\begin{example}
求$n$阶方阵$\begin{pmatrix}
x&     1&\cdots& 1\\
1&     x&\ddots&\vdots\\
\vdots&\ddots&     x& 1\\
1&\cdots&     1& x\\
\end{pmatrix}$的秩.
\end{example}
\pp
\begin{example}
求$n$阶方阵$\begin{pmatrix}
1&     1&      &     \\
 &     1&\ddots&     \\
 &      &\ddots&    1\\
1&      &      &    1\\
\end{pmatrix}$的秩.
\end{example}
\end{frame}






\begin{frame}\frametitle{例题}
\begin{example}
每个非零矩阵都为一个列满秩矩阵和一个行满秩矩阵的乘积.
\end{example}
\pp

\begin{example}
每个秩为$r$的矩阵都可以写成$r$个秩为$1$的矩阵的和.
\end{example}
\pp

\begin{example}
若$A$为列满秩$m\times n$矩阵,则$A$为某个$m$阶可逆阵的前$n$列.\\
若$A$为行满秩$m\times n$矩阵,则$A$为某个$n$阶可逆阵的前$m$行.
\end{example}
\end{frame}






\begin{frame}\frametitle{例题}
\begin{example}\frametitle{例题}
设矩阵$A$的列数和矩阵$B$的行数为$n$,则
\[\rank(A)+\rank(B)-n\leq \rank(AB)\leq \min\Big(\rank(A),\rank(B)\Big).\]
\end{example}
\pp

\begin{example}若$A$为$n$阶方阵满足$A^2=A$,则$\rank(A)+\rank(I-A)=n$.
\end{example}
\pp
\begin{example}
设$A\in F^{m\times n}$, $B\in F^{n\times m}$. 则
\[m+\rank(I_n-BA) = \rank\begin{pmatrix}I_m&A\\B&I_n\end{pmatrix} =n+\rank(I_m-AB).\]
\end{example}
\end{frame}







